% !TeX program = xelatex
\documentclass[12pt,a4paper]{article}
\usepackage{amsmath,amssymb}
\usepackage{geometry}
\usepackage{booktabs}
\usepackage{enumitem}
\usepackage[UTF8]{ctex}
\geometry{margin=2.5cm}
\title{\textbf{Steven Strogatz《非线性动力学与混沌》\\[4pt] 第2章与第3章\quad 中文翻译}}
\author{基于 OCR 文本提取 + 人工校订数学符号}
\date{}

\begin{document}
\maketitle

% ============================================================
\section{第2章\quad 直线上的流（Flows on the Line）}

\subsection{2.0\quad 引言}

第1章通过几个简单例子概述了非线性动力学的核心思想：几何思维、不动点、稳定性、
分岔，以及对混沌的简要预告。第2章将这些思想系统化。

我们将一阶自治系统写成：
\begin{equation}
\dot{x} = f(x)
\end{equation}
其中 $x(t)$ 是实值函数，$f(x)$ 是 $x$ 的光滑实值函数。这样的方程被称为一维的或一阶的，
因为只有一个变量 $x$ 出现，且方程只含 $x$ 的一阶导数。也称作\textbf{自治的}，因为 $f$ 不显含
时间 $t$。

\subsection{2.1\quad 几何思维方式}

\textbf{核心思想}：将 $\dot{x} = f(x)$ 解释为实线上的向量场。

想象一种流体（``相流体''）沿实线流动，其局部速度等于 $f(x)$。流速向右处 $f(x) > 0$，
向左处 $f(x) < 0$。为求得从任意初始条件 $x_0$ 出发的解，我们在 $x_0$ 处放置一个假想
粒子（``相点''），观察它如何被流体携带。随时间推移，相点沿实线按照某个函数 $x(t)$
运动，该函数称为基于 $x_0$ 的\textbf{轨迹}，代表了从初始条件 $x_0$ 出发的微分方程的解。

\textbf{不动点}（fixed points）$x^*$ 定义为 $f(x^*) = 0$——它们是流的驻点。实心黑点表示稳定
不动点（局部流向它），空心点表示不稳定不动点（流远离它）。

不动点代表\textbf{平衡解}（也称稳态解、常数解或静止解）。如果所有足够小的扰动在时间中
衰减，则称该平衡是稳定的。

\begin{quote}
\textbf{例 2.1.1}：分析 $\dot{x} = \sin x$。

\textbf{解}：不动点 $x^* = k\pi$（$k$ 为整数）。奇数倍 $\pi$ 是稳定不动点（流向它们），
偶数倍 $\pi$ 是不稳定不动点（流远离它们）。
\end{quote}

\subsection{2.2\quad 不动点与稳定性}

将上一节的思想推广到任意一维系统 $\dot{x} = f(x)$。我们只需要画出 $f(x)$ 的图像，然后用
它来勾画实线上的向量场。

\subsubsection{关键定义}
\begin{itemize}[nosep]
  \item \textbf{相流体}：以局部速度 $f(x)$ 沿实线流动的假想流体
  \item \textbf{相空间}：实线本身
  \item \textbf{相点}：被流携带的假想粒子
  \item \textbf{轨迹}：代表从给定初值出发的解的函数 $x(t)$
  \item \textbf{相图}：显示系统所有定性不同轨迹的图形
  \item \textbf{不动点} $x^*$：$f(x^*) = 0$ 的点——流的停滞点
  \item \textbf{稳定不动点}（实心黑点）：局部流向它
  \item \textbf{不稳定不动点}（空心点）：流远离它
  \item \textbf{稳定平衡}：所有足够小的扰动在时间中衰减
  \item \textbf{局部稳定 vs 全局稳定}：不动点可能对小扰动是吸引的，但对大扰动不是
\end{itemize}

\begin{quote}
\textbf{例 2.2.1}：求 $\dot{x} = x^2 - 1$ 的所有不动点并判断稳定性。

\textbf{解}：不动点 $x^* = \pm 1$。$x^* = -1$ 稳定，$x^* = 1$ 不稳定。

\medskip
\textbf{例 2.2.2}：RC 电路。$\dot{Q} = V_0/R - Q/(RC)$。不动点 $Q^* = CV_0$ 全局稳定。

\medskip
\textbf{例 2.2.3}：$\dot{x} = x - \cos x$。通过画 $y = x$ 和 $y = \cos x$ 找交点——无需不动点的显式公式
即可判断：唯一的不动点是不稳定的。
\end{quote}

\subsection{2.3\quad 种群增长}

指数增长模型：$\dot{N} = rN$，解为 $N(t) = N_0 e^{rt}$。这个模型在长期内不现实。

\textbf{逻辑斯谛（Logistic）方程}：
\begin{equation}
\dot{N} = rN\left(1 - \frac{N}{K}\right)
\end{equation}
由 Verhulst 于 1838 年首次提出。$r$ 为增长率，$K$ 为环境容纳量（carrying capacity）。
人均增长率随 $N$ 线性下降（图 2.3.2）。

不动点：$N^* = 0$（不稳定），$N^* = K$（稳定）。
从任意 $N_0 > 0$ 出发的相点总是流向 $N = K$。

\textbf{解的定性形状}：
\begin{itemize}[nosep]
  \item $N_0 < K/2$：先加速增长（上凸），经过 $N = K/2$ 后减速（下凸）$\to$ S 形曲线
  \item $K/2 < N_0 < K$：从开始就减速（一直下凸）
  \item $N_0 > K$：单调递减至 $K$（上凸）
\end{itemize}

\textbf{对逻辑斯谛模型的批评}：模型的代数形式不应从字面上理解。它应被视为一个隐喻——
描述从零增长到某个容纳量 $K$ 的趋势。对细菌和酵母等简单生物的实验室实验经常
给出令人印象深刻的匹配，但对具有复杂生命周期（卵、幼虫、蛹、成虫）的生物
（果蝇、面粉甲虫）则效果不佳。

\subsection{2.4\quad 线性稳定性分析}

设 $x^*$ 是一个不动点，令 $\eta(t) = x(t) - x^*$ 是一个小扰动。则：
\begin{equation}
\dot{\eta} = \dot{x} = f(x) = f(x^* + \eta)
\end{equation}

\textbf{泰勒展开}：
\begin{equation}
f(x^* + \eta) = f(x^*) + f'(x^*)\eta + O(\eta^2)
\end{equation}

因为 $f(x^*) = 0$（不动点定义），忽略 $O(\eta^2)$ 项得到\textbf{线性化}：
\begin{equation}\label{eq:linearization}
\dot{\eta} \approx f'(x^*)\,\eta
\end{equation}
这是关于 $\eta$ 的线性微分方程，称为关于 $x^*$ 的线性化。
\begin{itemize}[nosep]
  \item 若 $f'(x^*) > 0$：$\eta(t)$ 指数增长 $\to$ \textbf{不稳定}
  \item 若 $f'(x^*) < 0$：$\eta(t)$ 指数衰减 $\to$ \textbf{渐近稳定}
  \item 若 $f'(x^*) = 0$：需要非线性分析（图解法）——不动点可能是稳定、不稳定、
        半稳定，或是整条不动点线
\end{itemize}

\textbf{特征时间尺度}（characteristic time scale）：$\tau = 1/|f'(x^*)|$，决定了 $x(t)$ 在 $x^*$
邻域内发生显著变化所需的时间。

\begin{quote}
\textbf{例 2.4.1}：$\dot{x} = \sin x$。$f'(x^*) = \cos(k\pi)$。$k$ 为偶数时 $f' = 1$（不稳定），
$k$ 为奇数时 $f' = -1$（稳定）。

\medskip
\textbf{例 2.4.2}：逻辑斯谛方程。$f'(0) = r$（不稳定），$f'(K) = -r$（稳定）。
特征时间尺度 $= 1/r$。

\medskip
\textbf{例 2.4.3}：$f'(x^*) = 0$ 时的稳定性——需具体情况具体分析，用图解法判断。
四种情况：
\begin{enumerate}[label=(\alph*),nosep]
  \item $\dot{x} = -x^3$（稳定）
  \item $\dot{x} = x^3$（不稳定）
  \item $\dot{x} = x^2$（半稳定——从左边吸引、从右边排斥）
  \item $\dot{x} = 0$（整条不动点线——扰动既不增长也不衰减）
\end{enumerate}
\end{quote}

\subsection{2.5\quad 存在性与唯一性}

\begin{quote}
\textbf{例 2.5.1}：$\dot{x} = x^{1/3}$，$x(0) = 0$ 有多重解：
\[
x(t) = 0 \quad\text{和}\quad x(t) = \left(\frac{2t}{3}\right)^{3/2}
\]
事实上存在无穷多个解。原因：$f'(0)$ 为无穷大。
\end{quote}

\textbf{存在唯一性定理}：考虑初值问题 $\dot{x} = f(x), x(0) = x_0$。若 $f(x)$ 和 $f'(x)$ 在实轴的
某个开区间 $R$ 上连续，且 $x_0$ 属于 $R$，则该初值问题在 $t = 0$ 附近的某个时间区间
$(-\tau, \tau)$ 上存在解 $x(t)$，且解是唯一的。

该定理\textbf{不}保证解对所有时间都存在。

\begin{quote}
\textbf{例 2.5.2}：$\dot{x} = 1 + x^2$，$x(0) = 0$。分离变量得 $x(t) = \tan t$。解仅在
$-\pi/2 < t < \pi/2$ 内存在——在有限时间内达到无穷大。这种现象称为
\textbf{blow-up（爆炸）}，在燃烧等逃逸过程中具有物理意义。
\end{quote}

\subsection{2.6\quad 振荡的不可能性}

一阶系统中不动点主导动力学。所有轨迹要么趋近不动点，要么发散到 $\pm\infty$。
相位点永远不会反转方向。因此：
\begin{itemize}[nosep]
  \item 趋近平衡总是单调的——不会出现超调或阻尼振荡
  \item 不存在无阻尼振荡
  \item $\dot{x} = f(x)$ 没有周期解
\end{itemize}

这些结果是\textbf{拓扑性质}的：在直线上单调流动，永远不会回到出发点。
（在圆周上，向量场可以出现周期解——参见第4章。）

\textbf{机械类比（过阻尼系统）}：牛顿定律 $m\ddot{x} + b\dot{x} = F(x)$。当阻尼远大于惯性
（$b\dot{x} \gg m\ddot{x}$）时，系统近似为 $\dot{x} = f(x)$，其中 $f(x) = F(x)/b$。
质量块被缓慢拖回平衡位置，不会振荡。

\subsection{2.7\quad 势函数}

势函数 $V(x)$ 定义为：
\begin{equation}
f(x) = -\frac{dV}{dx}
\end{equation}

时间导数：
\[
\frac{dV}{dt} = \frac{dV}{dx} \cdot \frac{dx}{dt} = -\left(\frac{dV}{dx}\right)^2 \leq 0
\]
因此 $V(t)$ 沿轨迹递减——粒子总是向更低的势能处运动。
\begin{itemize}[nosep]
  \item $V(x)$ 的局部极小值 $\to$ 稳定不动点
  \item $V(x)$ 的局部极大值 $\to$ 不稳定不动点
  \item 势函数仅定义到差一个加法常数
\end{itemize}

\begin{quote}
\textbf{例 2.7.1}：$\dot{x} = -x$。势 $V(x) = x^2/2$。唯一的平衡点在 $x = 0$，稳定。

\medskip
\textbf{例 2.7.2}：$\dot{x} = x - x^3$。势 $V(x) = -\frac{x^2}{2} + \frac{x^4}{4}$。
称为\textbf{双阱势}（double-well potential），系统是\textbf{双稳态}的（bistable）：
$x = \pm 1$ 稳定，$x = 0$ 不稳定。
\end{quote}

\subsection{2.8\quad 在计算机上求解方程}

三种工具：解析方法、图解方法、数值方法。

\subsubsection{欧拉法（Euler's Method）}
一阶方法，误差 $\propto \Delta t$：
\begin{equation}
x_{n+1} = x_n + f(x_n)\Delta t
\end{equation}

\subsubsection{改进欧拉法（Improved Euler）}
二阶方法，误差 $\propto (\Delta t)^2$：
\begin{align}
\tilde{x}_{n+1} &= x_n + f(x_n)\Delta t \quad\text{（试探步）} \\
x_{n+1} &= x_n + \frac{1}{2}\bigl[f(x_n) + f(\tilde{x}_{n+1})\bigr]\Delta t \quad\text{（正式步）}
\end{align}

\subsubsection{四阶龙格-库塔法（Runge-Kutta）}
计算四个中间量并组合，在实际应用中达到准确性与计算成本的较好平衡。

\textbf{舍入误差}：计算机精度有限（单精度约 $10^{-7}$，双精度约 $10^{-16}$）。
步长过小会导致舍入误差累积。

\textbf{推荐软件}：MATLAB、Mathematica、Maple、Python、Sage、Slopes（入门）、
XPPAUT（研究级）。

% ============================================================
\section{第3章\quad 分岔（Bifurcations）}

\subsection{3.0\quad 引言}

第2章已经看到，直线上向量场的动力学非常有限：所有解要么趋近平衡，要么发散到
$\pm\infty$。给定动力学的这种平凡性，一维系统有什么有趣的呢？

\textbf{答案}：对参数的依赖。当参数变化时，流的定性结构可以改变。特别地，不动点可以被
创造或毁灭，或者其稳定性可以改变。这些动力学的定性变化称为\textbf{分岔}（bifurcations），
发生分岔的参数值称为\textbf{分岔点}（bifurcation points）。

分岔在科学上很重要——它们为控制参数变化时的相变和不稳定性提供了模型。
例如梁的屈曲：小重量时梁保持竖直（稳定）；重量过大时竖直位置变为不稳定，
梁可能弯曲。这里重量是控制参数，梁的偏转量是动力学变量 $x$。

\subsection{3.1\quad 鞍-结分岔（Saddle-Node Bifurcation）}

鞍-结分岔是不动点被创造和毁灭的基本机制。当参数变化时，两个不动点互相靠近、
碰撞并互相湮灭。

\subsubsection{标准型（Prototypical Example）}
\begin{equation}
\dot{x} = r + x^2
\end{equation}

\begin{itemize}[nosep]
  \item 当 $r < 0$：两个不动点（一个稳定、一个不稳定）
  \item 当 $r = 0$：一个半稳定不动点（$x^* = 0$）
  \item 当 $r > 0$：\textbf{没有不动点}
\end{itemize}

等价形式 $\dot{x} = r - x^2$：不动点为 $x^* = \pm\sqrt{r}$（$r > 0$ 时）。
$f'(x^*) = -2x^*$，所以 $x^* = +\sqrt{r}$ 稳定，$x^* = -\sqrt{r}$ 不稳定。
在 $r = 0$ 处，$f'(x^*) = 0$——不动点合并时线性化消失。

\textbf{图示惯例}：分岔图以 $r$ 为横轴、$x^*$ 为纵轴。实线 = 稳定，虚线 = 不稳定。

\textbf{术语}：也称为折叠分岔（fold bifurcation）、转折点分岔（turning-point bifurcation）、
蓝天分岔（blue sky bifurcation）——一对不动点``从晴朗的蓝天中''出现。

\subsubsection{正规形（Normal Forms）}
在鞍-结分岔附近，动力学通常形如 $\dot{x} = r \pm x^2$。
对 $f(x,r)$ 在分岔点 $(x^*, r_c)$ 处做泰勒展开：
\begin{equation}
\dot{x} \approx a(r - r_c) + b(x - x^*)^2
\end{equation}
其中 $a = \partial f/\partial r|_{(x^*,r_c)}$，$b = \frac{1}{2} \partial^2 f/\partial x^2|_{(x^*,r_c)}$
（假设 $a, b \neq 0$）。

\begin{quote}
\textbf{例 3.1.2}：$\dot{x} = r - x - e^{-x}$。不动点无法显式求解——用几何方法：同时画
$r - x$ 和 $e^{-x}$，看交点。分岔出现在两条曲线相切时：
\[
e^{-x} = r - x \quad\text{且}\quad -e^{-x} = -1
\]
解得 $x = 0, r = 1$。所以 $r_c = 1$。
\end{quote}

\subsection{3.2\quad 跨临界分岔（Transcritical Bifurcation）}

适用于某些科学情境：一个不动点必须对所有参数值都存在且永远不会被毁灭，
但可能改变稳定性。例如逻辑斯谛方程中零种群总是一个不动点。

\subsubsection{标准型}
\begin{equation}
\dot{x} = rx - x^2
\end{equation}

$x^* = 0$ 对所有 $r$ 都是不动点。
\begin{itemize}[nosep]
  \item $r < 0$：$x^* = 0$ 稳定，$x^* = r$ 不稳定
  \item $r = 0$：两个不动点合并
  \item $r > 0$：$x^* = 0$ 不稳定，$x^* = r$ 稳定
\end{itemize}
发生了两个不动点之间的\textbf{``稳定性交换''}（exchange of stabilities）。

\textbf{与鞍-结分岔的关键区别}：在跨临界情形中，两个不动点在分岔后不会消失——
它们只是交换了稳定性。

\begin{quote}
\textbf{例 3.2.1}：$\dot{x} = x(1 - x^2) - a(1 - e^{-bx})$。$x = 0$ 对所有 $(a,b)$ 都是不动点。
级数展开得：$\dot{x} \approx (1 - ab)x + \frac{ab^2}{2}x^2$。当 $ab = 1$ 时发生跨临界分岔。

\medskip
\textbf{例 3.2.2}：$\dot{x} = r\ln x + x - 1$ 在 $x = 1$ 附近。令 $u = x - 1$，展开：
$\dot{u} \approx (r + 1)u - \frac{r}{2}u^2$。在 $r_c = -1$ 处发生跨临界分岔。
通过尺度变换可化为正规形 $\dot{X} \approx RX - X^2$。
\end{quote}

\subsection{3.3\quad 激光阈值}

将分岔理论应用于固体激光器的简化模型（据 Haken 1983）。

\textbf{物理背景}：激光活性原子嵌入固态基质中，两端有部分反射镜。外部能量源``泵浦''
原子脱离基态。泵浦较弱时，原子独立振荡（普通灯）。泵浦强度超过阈值时，
原子突然开始同相振荡——灯变成了激光。这是自组织的合作现象。

\textbf{模型}（简化，忽略量子力学）：
\begin{itemize}[nosep]
  \item $n$ = 激光场中的光子数
  \item $\dot{n} = \text{gain} - \text{loss} = GnN - kn$
  \item $G$ = 增益系数，$k$ = 损耗率
  \item $N(t) = N_0 - \alpha n$（受激原子因发射光子而减少）
  \item $N_0$ = 无激光作用时的泵浦强度
\end{itemize}

代入得：
\begin{equation}
\dot{n} = (GN_0 - k)n - (\alpha G)n^2
\end{equation}

系统在 $N_0 = k/G$ 处经历\textbf{跨临界分岔}：
\begin{itemize}[nosep]
  \item $N_0 < k/G$：$n^* = 0$ 稳定（灯状态）
  \item $N_0 > k/G$：$n^* = 0$ 变为不稳定，出现新的稳定不动点
        $n^* = (GN_0 - k)/(\alpha G) > 0$（激光状态）
\end{itemize}

虽然该模型正确预测了阈值的存在，但它忽略了激发态原子的动力学、自发辐射等。

\subsection{3.4\quad 叉形分岔（Pitchfork Bifurcation）}

常见于具有对称性的物理问题（如左右空间对称性）。不动点以对称对的方式出现和消失。

\subsubsection{超临界叉形分岔（Supercritical Pitchfork Bifurcation）}

标准型：
\begin{equation}
\dot{x} = rx - x^3
\end{equation}

该方程在 $x \to -x$ 变换下不变——这是左右对称性的数学表达。

\begin{itemize}[nosep]
  \item $r < 0$：原点是唯一的不动点，稳定
  \item $r = 0$：原点仍然稳定但更弱——线性化消失。解不再指数衰减，而是以代数函数
        形式缓慢衰减。这种惰性衰减称为\textbf{临界慢化}（critical slowing down）
  \item $r > 0$：原点变为不稳定，两个新的稳定不动点对称地出现在 $x^* = \pm\sqrt{r}$
\end{itemize}

分岔图的形状类似``干草叉''（pitchfork）。

\begin{quote}
\textbf{例 3.4.1}：$\dot{x} = -x + \beta\tanh x$。方程出现在磁体和神经网络的统计力学模型中。
$\tanh x$ 在原点处的斜率为 $\beta$。
\begin{itemize}[nosep]
  \item $\beta < 1$：唯一不动点（原点，稳定）
  \item $\beta = 1$：在 $x^* = 0$ 处发生叉形分岔
  \item $\beta > 1$：原点变不稳定，出现两个新的稳定不动点
\end{itemize}
\end{quote}

% ============================================================
\vspace{12pt}
\hrule
\vspace{12pt}

\section*{翻译说明}

\begin{enumerate}[nosep]
  \item 本翻译基于 PDF OCR 文本提取（Tesseract），数学符号已经人工校订修正。
  \item OCR 文件在第3章\S3.4（超临界叉形分岔）处截断。缺失的后续章节包括：
        \begin{itemize}[nosep]
          \item 亚临界叉形分岔（标准型 $\dot{x} = rx + x^3$，含跳跃现象和滞后）
          \item 旋转圆环上的过阻尼珠子（连接分岔的力学类比）
          \item 不完全分岔与突变理论
          \item 虫害爆发模型（云杉蚜虫）与尖点突变
          \item 第3章习题
        \end{itemize}
  \item 页码引用（如``第26页''）指原版教材页码，非 PDF 页码。
\end{enumerate}

\end{document}
